\chapter{Краткое описание модели с линейными порогами}

Модель с линейными порогами (Linear Threshold Model, LTM) --- это модель распространения влияния в графах, где каждому узлу назначается порог активации. Узел активируется, если доля его активных соседей превышает назначенный порог. Для выбора k наиболее влиятельных узлов применяется жадный алгоритм, максимизирующий ожидаемое количество активированных узлов. Процесс проводится итерационно до достижения состояния равновесия. Модель оценивает способность узлов вызывать каскадные активации и может использоваться для анализа влияния внутри выявленных сообществ.



\clearpage
